% Section 7: M-SIGN (PHASE5a)
\section{The sign module M-SIGN: R3, R2 and R10}\label{sec:msign}

Source: PHASE5a (Appendix~\ref{app:phase5a}), proof-register items R3, R2 and
R10. Checks: \texttt{verify\_R3.py}, \texttt{verify\_R2.py},
\texttt{verify\_R10.py}. \textbf{Scope: every statement is for the pooled LCE
model M0.} The sign and CR rows are valid linear rows under the district rungs
too (Prop.~M5(a)), but sharpness and inference there are proof item R16, not
proved. The package labels the module \emph{experimental} in 0.3.0 and 0.3.1.

\paragraph{Setting.} $\ThI$ and $\ThO$ are viewed in the free coordinates $k$.
The target is a model-implied linear $\varphi(k)=\varphi_0+h^\top k$ (for example
$\tilde B_{rc}$ or the gap $\tilde D=\tilde B_1-\tilde B_2$). An \emph{admissible
restriction} is a set $C(P)$ that is (C1) closed and convex and (C2) depends on
$P$ only through the design and $Q(P)$. Examples: any sign on $\delta^*_B$,
$\Delta^*_r$ or $\delta^{r*}_W$; mean-level ordering on $A$; the M-LOC rows. A
\emph{union restriction} is $C_+\cup C_-$; contextual reinforcement (CR) for a pair
$(g_B,g_W)$ is $C_\pm=\{\pm g_B\ge0,\pm g_W\ge0\}$.

\subsection{R3: sharpness of LCE plus restrictions}

\begin{result}[Theorem R3 --- PHASE5a \S1]\label{res:r3}
\status{Proved under conditions} (pooled M0). Assume A1--A5, standing
integrability, $P(X\in A)=1$ and $\operatorname{affdim}(A)=R-1$.
\begin{enumerate}[label=(\alph*),nosep]
\item For admissible $C$, the sharp identified set for $\varphi$ under LCE $\cap\,C$
is $[\min_{\ThI\cap C}\varphi,\max_{\ThI\cap C}\varphi]$ or empty.
\item $\ThO\cap C=\bigcup_{P'\in\mathcal P(P)}(\ThI(P')\cap C(P'))$: the LP
interval is sharp for the regression and the design; a finite LP when $C$ is
polyhedral; equal to (a) when $\ThI=\ThO$; weakly outer otherwise.
\item For a union, the set is $\varphi(\Theta\cap C_+)\cup\varphi(\Theta\cap C_-)$,
at most two intervals; the hull is sharp iff they intersect; a common witness in
$\Theta\cap C_+\cap C_-$ certifies that (sufficient only).
\item One-vs-rest (W1) CR is one LP: $C_+\cup C_-=\{D^{r*}(k)\in\operatorname{hull}\{\gamma_rER^{r*},ER^{r*}\}\}$,
two linear rows; in $2\times2$ this is the only CR, and under it
$\operatorname{sign}\tilde D=\operatorname{sign}ER^*$.
\item Refutation is joint: an empty population LP rejects the conjunction LCE
$\cap\,C$, not either part alone.
\end{enumerate}
\end{result}
\begin{plain}
Adding a sign or ordering belief to the pooled linear rung keeps the bound
``as tight as the regression allows'': nothing narrower can be learned from the
regression and the design. The full data could rule out more (Remark~R3.1). For
the default one-vs-rest comparison, reinforcement becomes a single band in the
gap, so no union is needed. If the bound is empty, it is the combination that is
contradicted, not the belief on its own.
\end{plain}
\begin{sketch}
(a) The identified set is the image of the convex compact $\ThI\cap C$ under an
affine map. (b) Apply the construction of PHASE2a Prop.~2 to each $k\in\ThO\cap
C$; (C2) keeps $C$ unchanged across laws with the same regression. (d) Apply
Lemma~M1(c) to the model means: $\delta^{r*}_W=\Var_p(X_r)(ER^{r*}-D^{r*})$ and
$\delta^{r*}_B=\mu_r(1-\mu_r)(D^{r*}-\gamma_rER^{r*})$; the union of the two
same-sign cases is the hull.
\end{sketch}
\proofref{Appendix~\ref{app:phase5a}, \S1.3, steps 1--15; checks S1--S7 (\S1.5)}
\begin{revnote}
(M-1; P-B15, P-B16, P-B17; I.7.) The dimension hypothesis is
$\operatorname{affdim}(A)=R-1$, not ``nonempty relative interior''. ``The LP
interval'' is a finite LP only for polyhedral $C$, and in round~2 the
finite-LP sentence was narrowed to $\ThO\cap C$. (e) earlier read ``never $C$
alone''; it now says the conjunction is rejected and gives both directions of
counterexample. The scope was restricted to pooled M0 (P-B7).
\end{revnote}

\paragraph{Remark R3.1 (full law versus regression).} On $X\in\{1/4,1/2,3/4\}$ with
a fair binary $T$ at every composition, the full law forces $k=0$, but the
regression admits a range. For $C=\{k\in[.2,.4]\}$ the full-law set is empty
while the regression-only gap set is $[1/6,1/3]$ (check S6 [V]).

\paragraph{Remark R3.2 (LCE can exclude CR although LCE fits).} On
$X\in\{0,1/4,1\}$ with $m_0=1-x$, $m_1=(49/55)x$, the vertex rows force $k=0$, so
LCE fits with a single point, $ER=0$ and $\tilde D=-261/1925$; both CR branches
are empty. In the model-free continuous problem CR is still feasible. Any ``cannot
be refuted'' statement must name its base system (check S7 [F]).

\paragraph{\S1.4: what fails.} For CR pairings other than W1, (d) fails: both
branches can be non-empty with disjoint images (search S5 found gaps up to about
31\% of hull width). For the pairwise version only $\delta^*_B(r,s)$ is pairwise;
$\Delta^*_r$ pools all groups (E2-1). N2 shapes, N1 realised signs and targets
mixing $\tilde B$ with $B$ are not covered.

\subsection{R2: when does an added restriction tighten the bound?}

\begin{result}[Proposition R2 --- PHASE5a \S2]\label{res:r2}
\status{Proved.} Let $F$ be the current (non-empty, compact) feasible set,
$F^*_U,F^*_L$ its optimal faces for $\varphi$, and $N$ an admissible restriction
with $F\cap N\ne\emptyset$.
\begin{enumerate}[label=(\alph*),nosep]
\item $U_N<U$ iff $F^*_U\cap N=\emptyset$ (one feasibility LP, carrying the
target intercept: $h^\top k\ge U-\varphi_0$).
\item For one row $a^\top k\le b$: tightening iff $\min_{F^*_U}a^\top k>b$ iff
every optimal dual multiplier is positive. A positive multiplier from one solver
optimum does not certify tightening.
\item $N$ is implied iff it tightens no linear target (the witness need not be a
reported target).
\item In $2\times2$ under LCE with $\gamma<1$, each of the four sign rows is either
implied or tightens any target non-constant in $w_1$; closed forms are those of G2.
\item For stacks, report added-first and added-last gains; zero added-last gain
means zero incremental gain for this target, not that the row is implied.
\end{enumerate}
\end{result}
\begin{plain}
A belief narrows a quantity only if it removes every table that attains that end
of the bound. The package tests this with one feasibility LP per end and never
trusts solver dual values. A belief can narrow the set of possible tables without
narrowing the quantity you care about; with three or more groups this was common
in simulation (28\% of cases).
\end{plain}
\proofref{Appendix~\ref{app:phase5a}, \S2.3, steps 1--14; checks T1--T4, T1b, T1c, T4b}
\begin{revnote}
(P-B18--P-B22.) The face LP in step~4 omitted the target intercept (the code
was correct; check T1c [F]); the duplicate-row example shows a positive dual does
not certify tightening; (d) gained its qualifications; (e)'s ``redundant given
the others'' read as ``implied'' and was reworded (T4b [F]); width-gain
inference must treat both endpoints jointly.
\end{revnote}

\subsection{R10: inference with added rows, unions and near-empty sets}

\begin{result}[Theorem R10 --- PHASE5a \S3]\label{res:r10}
\status{Proved under conditions} (pooled M0; by interval version, below). Assume
P3.5's conditions for the augmented row system and that the restrictions hold at
the truth.
\begin{enumerate}[label=(\alph*),nosep]
\item \emph{Added rows.} For finitely many rows linear in $\beta$, P3.5 and
P3.6(a) hold for the augmented system, version by version. For actual contrasts
the residual allowance must cover the contrast (a contrast envelope, or the
per-cell sum the package uses).
\item \emph{Parallel disjunction (version (a) only).} Inflate each piece, then take
the maximum; this covers at the nominal level. W1 and $2\times2$ CR are of this
form (the ``slab''); no branch is selected from $\widehat{ER}$.
\item \emph{General union.} $\mathrm{CI}(C_+)\cup\mathrm{CI}(C_-)$, each at the
full level, covers the true target with probability $\ge1-\alpha$ asymptotically;
no Bonferroni split is needed. This is point coverage only.
\item \emph{Near-empty sets.} $P(\mathrm{CI}=\emptyset)\to\le\alpha_U+\alpha_L$
under the model, including point-identified cases; conditioning on non-emptiness
can only raise coverage. If the restrictions are false, no coverage statement
holds; a narrow non-empty interval is no evidence that they are true.
\item \emph{Plug-in brackets are not valid.} Imposing the estimated bracket
$\operatorname{hull}\{\gamma_r\widehat{ER},\widehat{ER}\}$ without inflation can
under-cover persistently: in a repeated three-point design
$\limsup_nP(\tilde D_0\in\mathrm{CI}_{\rm naive})\le\Phi(\tfrac7{12}z_{.995})=0.9335<0.95$.
\end{enumerate}
\end{result}

\begin{longtable}{@{}L{0.30\linewidth}L{0.42\linewidth}L{0.20\linewidth}@{}}
\caption{Status of R10 by interval version (PHASE5a \S3.2).}\label{tab:r10}\\
\toprule Interval version & What R10 claims & Status \\ \midrule\endfirsthead
\toprule Interval version & What R10 claims & Status \\ \midrule\endhead
P3.5(a), conservative & (a)--(d), including the slab (b) & Proved \\
P3.5(b1), CLR over all dual vertices & (a), (c), (d) under the full augmented CLR conditions & Cited \\
P3.5(b2)(i), fixed superset & branch union and crossing arguments & Proved \\
P3.5(b2)(ii), projection & both ends, any candidate set; $\Sigma$ p.d.\ or exact reduction & Proved \\
P3.5(b3), CLR on simulated candidates (package default for CR) & the inclusion and crossing steps are deterministic; their probabilities are conjectural & Algorithm \\
\bottomrule
\end{longtable}

\begin{plain}
Sign beliefs enter the same confidence-interval machinery as the rung's own
constraints. Reinforcement uses a band whose two pieces are each widened by their
own uncertainty before the larger is taken; this is valid without knowing which
branch is true. The tempting shortcut, plugging the estimated neighbourhood and
Goodman numbers into the bracket, is not a confidence interval, and more data do
not cure it. When the belief is false, the interval stays narrow and can be
reproducibly wrong (Section~\ref{sec:evidence}).
\end{plain}
\proofref{Appendix~\ref{app:phase5a}, \S3.3, steps 1--23 (step 18--23 is the referee's fixed-design proof of (e)); checks V9--V13}
\begin{revnote}
(M-2; P-B23--P-B27, P-E17, P-E21, P-E22, P-E33; I.2; E2-2, E2-3.) The status of
each claim is now given by interval version (table above) with CLR's
Condition~P spelled out; zero-variance candidates are exact only with zero
gradient; the contrast envelope was added; (b)--(d) gained the one-event,
point-coverage and triangular-array qualifications and the belief-selection
sentence; (e) became an analytic theorem (referee's proof, credited) instead of a
Monte Carlo sketch; ``a fortiori'' was replaced by a direct $\le\tfrac12$ argument
(V10 [F]); union lengths are occupied lengths (V11). In round~2 the R10 scope was
narrowed to pooled M0 (I.2), and the coverage figure was tied to the procedure:
the literal intersection covered 304/600, CR-naive 550/600 (V13 [V]).
\end{revnote}

\paragraph{Consequences for the package (PHASE5a \S\S1.6, 2.5, 3.5).} For pooled
M0 with sign, ordering or local rows the LP interval may be labelled ``sharp for
the regression''; under district rungs, ``outer for the regression''. W1 and
$2\times2$ CR use the slab for the conservative interval and the branch union for
CLR (the default, algorithmic). Other CR pairings use the union of branch CIs and
can print two or more pieces. Redundancy verdicts use the face test. Emptiness is
reported as a test result, naming the rejected combination and whether the rung
alone is rejected.
